\section{Introduction} \label{sec:introduction}
  % Setting up the scene and protagonist [DRAM]
    In modern computing systems Memory serves as a fundamental part, directly influencing overall 
    performance, energy efficiency, and system reliability~\cite{paperRHARetrospective}. Ensuring that stored information 
    cannot be altered without authorization is a basic security requirement, as unauthorized 
    modifications would compromise both the integrity of computations and the system as a whole~\cite{paperTriggeringRHHardwareFault}.
    \\
    In \gls{dram}, each bit of data is stored in a single memory cell composed of a capacitor and a 
    transistor~\cite{paperRowhammeronaRPi}. Over decades of technological progress, improvements in transistor and capacitor 
    scaling have enabled \gls{dram} to serve as the predominant technology used for main memory across a wide 
    range of devices, from cloud servers and laptops to mobile phones~\cite{paperOnDRAM}.
    \\
    The demand for increasingly large memory capacities within a limited physical footprint has driven 
    continuous reductions in \gls{dram} cell size and increases in density. These changes allow memory cells 
    to operate with lower voltages and smaller charges, resulting in higher storage capacity and reduced 
    energy consumption~\cite{paperAnotherFlip}.
    \\
    The extensive use of \gls{dram} as main memory implies that potential security flaws within the technology 
    can have broad impact on dependent computing systems~\cite{paperOnDRAM}. Operations performed on a given memory 
    address must remain strictly confined to that address and must not inadvertently affect the contents 
    of other memory locations~\cite{paperRHARetrospective}. As memory technologies continue to scale toward smaller physical 
    dimensions, \gls{dram} chips become increasingly susceptible to disturbance effects, in which the behavior 
    of a memory cell can unintentionally influence adjacent cells~\cite{paperRHARetrospective}, causing the stored value of a 
    \gls{dram} cell to change without explicit modification, a phenomenon known as a bitflip~\cite{paperAnotherFlip}.

  % Conflict [Rowhammer]
    In 2014, Kim et al.~\cite{paperFlippingBits} demonstrated that \gls{dram} is vulnerable to disturbance errors, showing that 
    repeated activations of a memory row can induce charge leakage in neighboring cells. Such leakage 
    may alter stored values, which in turn can be exploited to escape sandboxes or escalate privileges 
    on a local system~\cite{paperRowhammeronaRPi}
    \\
    Initial verification of the Rowhammer phenomenon was conducted using FPGA-based memory controllers. 
    Later studies extended these attacks to native x86 systems as well as ARM-based platforms~\cite{paperRowhammeronaRPi}. 
    Early speculation suggested that Rowhammer attacks on ARM devices might be infeasible, largely due 
    to the presumed slower response of ARM memory controllers. Nonetheless, this assumption was later disproven 
    when bitflips were successfully observed on ARM hardware~\cite{paperDrammer}.

  % Research Question
    Previous studies primarily focused on mobile phones or Chromebooks~\cite{paperHalfDouble, paperDrammer, paperDRAMA, paperBlacksmith}. Subsequently 
    on older Raspberry Pi models, leaving it uncertain whether the latest generations exhibit similar 
    vulnerabilities. This uncertainty is further compounded by the differences in \gls{lpddr} memory, which 
    due to power and die area constraints, behaves differently from standard \gls{ddr} memory~\cite{paperBlacksmith}. In the 
    context of this Bachelor thesis, we want to analyze to what extent modern Raspberry Pi platforms, 
    specifically the Raspberry Pi 4B and Raspberry Pi 5, are susceptible to bitflips when subjected to 
    Rowhammer access patterns.
    \\

  % Outline
    \noindent
    \textbf{Outline.} The thesis is organized into six sections. 
    \Cref{sec:background} "\textit{Background}" introduces cache and \gls{dram} architectures and general information of the Rowhammer bug and techniques. 
    \Cref{sec:methodology} "\textit{Methodology}" describes the experimental design, measurement procedures, and evaluation criteria. 
    \Cref{sec:overview-of-challenges} "\textit{Overview of Challenges}" describes the fundamental techniques required to execute Rowhammer attacks effectively, including uncached, targeted, and fast access patterns by analyzing and evaluating the corresponding experiments of these primitives. 
    \Cref{sec:uncached} "\textit{Direct DRAM Activation}", \cref{sec:targeted} "\textit{Address Mapping Reconstruction}", and \cref{sec:fast} "\textit{Outpacing DRAM Refresh}" build on this overview and examine these techniques in greater detail.
    \Cref{sec:related-work} "\textit{Related Work}" reviews prior research on Rowhammer attacks. 
    \Cref{sec:discussion-and-limitations} "\textit{Discussion and Limitations}" examines the implications of the results and the study's constraints. 
    \Cref{sec:conclusion} "\textit{Conclusion}" summarizes the findings and outlines potential directions for future work.
